% =========================================================
% Soundness and Completeness of First-Order Logic
% =========================================================

\subsection{Proof Theory: Soundness and Completeness}

% ---------------------------------------------------------
% TOOLKIT
% ---------------------------------------------------------

% ---------------------------------------------------------
% Soundness and Completeness (Definitions)
% ---------------------------------------------------------
\begin{definitionbox}{Definition (Soundness)}
\begin{definition}[Soundness]\label{def:fol-soundness}
A proof system is \emph{sound} if every provable formula is valid:
\[
\Gamma \vdash \varphi \implies \Gamma \models \varphi.
\]
\end{definition}
\end{definitionbox}
\begin{remark*}[Standard quantified statement]
For every admissible object satisfying the hypotheses of the statement, the
assertion holds in the formal language under discussion.
\end{remark*}

\begin{remark*}[Interpretation]
The statement records the named logical notion or result together with its
intended syntactic or semantic role.
\end{remark*}

\begin{dependencies}
\begin{itemize}
  \item \hyperref[def:fol-satisfaction]{Satisfaction}
  \item \hyperref[def:log-consequence]{Logical Consequence}
\end{itemize}
\end{dependencies}

\begin{definitionbox}{Definition (Completeness)}
\begin{definition}[Completeness]\label{def:fol-completeness}
A proof system is \emph{complete} if every valid formula is provable:
\[
\Gamma \models \varphi \implies \Gamma \vdash \varphi.
\]
\end{definition}
\end{definitionbox}
\begin{remark*}[Standard quantified statement]
For every admissible object satisfying the hypotheses of the statement, the
assertion holds in the formal language under discussion.
\end{remark*}

\begin{remark*}[Interpretation]
The statement records the named logical notion or result together with its
intended syntactic or semantic role.
\end{remark*}

% ---------------------------------------------------------
% Soundness Theorem for FOL
% ---------------------------------------------------------

\begin{dependencies}
\begin{itemize}
  \item \hyperref[def:fol-satisfaction]{Satisfaction}
  \item \hyperref[def:log-consequence]{Logical Consequence}
\end{itemize}
\end{dependencies}

\begin{theorembox}{Theorem (Soundness of First-Order Logic)}
\begin{theorem}[Soundness of First-Order Logic]\label{thm:fol-sound}
The standard inference rules for first-order logic (UI, UG, EI, EG, and the
propositional rules) are sound: they preserve truth in all structures.

If $\Gamma \vdash \varphi$, then $\Gamma \models \varphi$.
\hyperref[prf:fol-sound]{Go to proof.}
\end{theorem}
\end{theorembox}
\begin{remark*}[Standard quantified statement]
$\forall \Gamma,\varphi:\; \Gamma\vdash\varphi \;\to\; \forall\mathcal{M},\forall s,\; (\forall\gamma\in\Gamma,\; \mathcal{M},s\models\gamma) \to \mathcal{M},s\models\varphi$.
\end{remark*}

\begin{remark*}[Predicate reading]
The predicate form records the relation or property named by the statement in
the displayed logical context.
\end{remark*}

\begin{remark*}[Interpretation]
The statement records the named logical notion or result together with its
intended syntactic or semantic role.
\end{remark*}

\begin{dependencies}
\begin{itemize}
  \item \hyperref[def:fol-soundness]{Soundness}
  \item \hyperref[def:fol-satisfaction]{Satisfaction}
  \item \hyperref[def:log-consequence]{Logical Consequence}
  \item \hyperref[ax:predicate-logic-universal-instantiation]{Universal Instantiation}
  \item \hyperref[ax:predicate-logic-quantifier-distribution]{Quantifier Distribution}
\end{itemize}
\end{dependencies}

% ---------------------------------------------------------
% G\"{o}del's Completeness Theorem
% ---------------------------------------------------------

\begin{theorembox}{Theorem (G\"{o}del's Completeness Theorem)}
\begin{theorem}[G\"{o}del's Completeness Theorem]
\label{thm:godel-complete}
First-order logic is complete: every logically valid formula is provable.
\[
\Gamma \models \varphi \implies \Gamma \vdash \varphi.
\]
Equivalently: if a set of sentences $\Gamma$ is consistent (has no proof of
contradiction), then $\Gamma$ has a model.
\hyperref[prf:godel-complete]{\textit{Go to proof.}}
\end{theorem}
\end{theorembox}
\begin{remark*}[Standard quantified statement]
For every admissible object satisfying the hypotheses of the statement, the
assertion holds in the formal language under discussion.
\end{remark*}

\begin{remark*}[Interpretation]
The statement records the named logical notion or result together with its
intended syntactic or semantic role.
\end{remark*}


\begin{remark*}[Exposition]
Soundness ensures proofs do not lead from true premises to false conclusions.
Completeness ensures proofs are powerful enough: every statement that is true
in all models can be established by a formal proof.

Together, soundness and completeness show that syntactic provability ($\vdash$)
and semantic entailment ($\models$) coincide for first-order logic.
\end{remark*}

\begin{remark*}[Exposition]
G\"{o}del's completeness theorem (1930) should not be confused with his
incompleteness theorems (1931). The completeness theorem concerns first-order
logic as a proof system. The incompleteness theorems concern the limitations of
formal theories capable of expressing arithmetic, and are a separate result.
\end{remark*}

\begin{dependencies}
\begin{itemize}
  \item \hyperref[def:fol-satisfaction]{Satisfaction}
  \item \hyperref[def:log-consequence]{Logical Consequence}
\end{itemize}
\end{dependencies}
